
\section{\trProjective{} $k$-\trScheme{}的\trSerreDuality}

凭借\trProjectiveSpace{}的\trSerreDuality，我们现在来证明\trProjective{} $k$-\trScheme{}的\trSerreDuality{}的种种版本。

\textbf{29.3.1. 回顾古怪的函子 $\pi^{!?}$。}

关键的构造是\trPushforward{} $\pi_*$ 的一个\emph{\trRightAdjoint} $\pi^{!?}$（当 $\pi$ 是一个\trFinite\trMorphism{}时），见 \S17.1.6 中的讨论。若你尚未读过那一节，现在应当读一读。

\textbf{29.3.2.}\ \emph{警示。}\quad 我们只对 $Y$ 上的\trQCSheaf{}定义过 $\pi^{!?}$，并非对一切 $\OO_Y$-\trModule。若 $\mathscr{G}$ 是一个 $\OO_X$-\trModule{}（不必\trQC），那么一般说来，如何为这个构造赋予意义、从而定义出一个 $\OO_Y$-\trModule{} $\pi^{!?}\mathscr{G}$，并\emph{不}清楚。（试试看！）然而，在 $\pi$ 是\trClosedEmbedding{}的特殊情形下，我们\emph{能够}理解 $\pi^{!?}\mathscr{G}$，如\mbox{练习~29.4.A} 所描述的那样。

\textbf{29.3.3.}\ \emph{注记。}\quad 若 $\pi$ \trFinite{}且\trFlat——这正是我们最感兴趣的情形——则 $\pi^{!?}$ 与 $\pi^!$ 相合（就「相合」一词唯一可能的含义而言），而后者是\trGrothendieck{}「六运算」（six operations）之一；我们的记号也正源于此（「$\pi^!$ 的一个奇怪变体」）。但 $\pi^!$ 天然地生活在\trDerivedCategory{}的世界里，所以我们在此不讨论它。（关于 $\pi^!$ 的更多内容，见 [KS1, Ch.~III]。）

现在我们把 $\pi^{!?}$ 的这套机制应用于\trSerreDuality。

\textbf{29.3.4. 射影 $n$ 维概形是 $\mathbb{P}^n$ 的\trFinite\trCover（finite cover）。}

\textbf{29.3.5. \trhProposition。}\ ——\ \emph{设 $X$ 是\trDimension{}至多为 $n$ 的\trProjective{} $k$-\trScheme。}

\emph{(a)~存在一个\trFinite\trMorphism{} $\pi\colon X\to\mathbb{P}^n$。}

\emph{(b)~若进而 $X$ 是\trPureDimension{}为 $n$ 的 \trCM，则 $\pi$ 是\trFlat{}的。}

\textbf{29.3.6.}\ \emph{致专家的注记。}\quad 我们在这里刻意允许 $n$ 大于 $\dim X$，这不是因为它有用，而是要指出：这个证明是行得通的，而一旦假设不当，我们证得的便是某种奇怪的东西。当 $\dim X=n$ 而 $X$ 并非 \trCM{}时，\mbox{命题~29.3.5(a)} 仍然有趣。而当 $\dim X<n$ 时，\mbox{命题~29.3.5(a)} 又蕴涵什么呢？

\emph{证明。}\quad (b) 部分由 (a) 部分与\trMiracleFlat{}定理~26.2.11 得出，所以我们来证 (a)。

选取一个\trClosedEmbedding{} $j\colon X\hookrightarrow\mathbb{P}^N$。为叙述简明，先设 $k$ 是无限域。由\mbox{练习~12.3.D(d)}，存在一个余维数为 $n+1$（比补维数恰小 $1$）且与 $X$ 不相交的\trLinearSpace{} $L$。从 $L$ 作\trProjection{}便给出一个态射 $\pi\colon X\to\mathbb{P}^n$。态射 $\pi$ 是\trAffine{}的（$\mathbb{P}^N\setminus L\to\mathbb{P}^n$ 是\trAffine{}的，而\trClosedEmbedding{} $X\hookrightarrow\mathbb{P}^N\setminus L$ 同一切\trClosedEmbedding{}一样是\trAffine{}的），又是\trProjective{}的（由射影态射的消去定理，\mbox{练习~17.3.C}；或者更直接地由\mbox{练习~17.3.D}），于是 $\pi$ 是\trFinite{}的（局部\trNoether{}概形上的射影\trAffine{}态射是\trFinite{}的，\mbox{推论~18.1.5}）。

\textbf{29.3.A.（练习）}\quad 在不假定 $k$ 无限的情形下证明\mbox{命题~29.3.5(a)}。提示：可先用\mbox{练习~12.3.D(c)} 证明：存在某个 $d$，使得有 $n+1$ 张 $d$ 次\trHypersurface{}的交与 $X$ 不相交。然后把上述论证应用于 $\mathbb{P}^N$ 的 $d$ 次\trVeronese{}嵌入（\S9.3.6）即可。\hfill$\square$

\textbf{29.3.7.}\ \emph{注记。}\quad 更一般地，上述论证表明：任何维数为 $n$（不必纯维）的\trProjective{} $k$-\trScheme{}都能表示为 $\mathbb{P}^n$ 的一个\trFinite\trCover。这或许可以看作\trNoetherNormalization{}的一个\trProjective{}版本。

\textbf{29.3.8. 借助 $\pi^{!?}$ 的 $X$ 上的\trSerreDuality。}

设 $\pi\colon X\to Y$ 是\trPureDimension{}为 $n$ 的\trProjective{} $k$-\trScheme{}之间的一个\trFinite\trMorphism。（我们马上就会在 $Y=\mathbb{P}^n_k$ 的情形下应用它，但我们不妨避开干扰，也无须那样特化。）

\textbf{29.3.9. \trhProposition。}\ ——\ \emph{设 $\pi$ \trFlat。若\trVectorBundle{}的函子版本\trSerreDuality{}对 $Y$ 成立，带\trCanonical{}层 $\omega_Y$，则\trVectorBundle{}的函子版本\trSerreDuality{}对 $X$ 成立，带\trCanonical{}层 $\pi^{!?}(\omega_Y)$。}

\emph{证明。}\quad 对每个 $i$ 以及 $X$ 上每个有限秩的\trLocallyFreeSheaf{} $\mathscr{F}$，我们有诸\trIsomorphism{}（随 $\mathscr{F}$ 函子）：
\[
{\small
\begin{aligned}
\mathrm{H}^{n-i}(X,\mathscr{F}^{\vee}\otimes\pi^{!?}\omega_Y)\cong{}& \mathrm{H}^{n-i}\bigl(Y,\pi_*(\mathscr{F}^{\vee}\otimes\pi^{!?}\omega_Y)\bigr) && \text{（$\pi$ 的\trAffineness，\S18.1~(v)）}\\
\cong{}& \mathrm{H}^{n-i}\bigl(Y,\pi_*(\mathscr{H}om_X(\mathscr{F},\pi^{!?}\omega_Y))\bigr) && \text{（\mbox{练习~14.1.G}）}\\
\cong{}& \mathrm{H}^{n-i}\bigl(Y,\mathscr{H}om_Y(\pi_*\mathscr{F},\omega_Y)\bigr) && \text{（\mbox{式~(17.1.8.2)}）}\\
\cong{}& \mathrm{H}^{n-i}\bigl(Y,(\pi_*\mathscr{F})^{\vee}\otimes\omega_Y\bigr) && \text{（\mbox{练习~24.3.E} 与\mbox{练习~14.1.G}）}\\
\cong{}& \mathrm{H}^{i}(Y,\pi_*\mathscr{F})^{\vee} && \text{（$Y$ 上\trVectorBundle{}的\trSerreDuality）}\\
\cong{}& \mathrm{H}^{i}(X,\mathscr{F})^{\vee}. && \text{（$\pi$ 的\trAffineness，\S18.1~(v)）}
\end{aligned}}
\]
（\trFlatness{}假设用于证明 $\pi_*\mathscr{F}$ 是\trLocallyFree{}的。）\hfill$\square$

\textbf{29.3.10. \trhCorollary。}\ ——\ \emph{\trVectorBundle{}的函子版本\trSerreDuality{}对一切纯维的 \trCM{} \trProjective{} $k$-\trScheme{}成立。}

\emph{证明。}\quad 把\mbox{命题~29.3.9} 与\mbox{命题~29.3.5(b)} 以及 \S29.2.2 结合起来。\hfill$\square$

\textbf{29.3.11. \trhProposition。}\ ——\ \emph{设 $\pi\colon X\to Y$ 是\trProjective{} $k$-\trScheme{}之间的一个\trFinite\trMorphism。若 $\mathrm{Hom}$ 的\trSerreDuality{}对 $Y$ 成立，带\trCanonical{}层 $(\omega_Y,t_Y)$，则 $\mathrm{Hom}$ 的\trSerreDuality{}对 $X$ 成立，带\trCanonical{}层 $\pi^{!?}\omega_Y$。}

回忆一下\mbox{练习~29.1.A}：对 $\mathrm{Hom}$ 的\trSerreDuality{}而言，迹版本与函子版本等价。因此我们径直使用「$\mathrm{Hom}$ 的\trSerreDuality」这个说法。

注意，我们对 $\pi$ 不作任何\trFlat{}假设（不同于\trVectorBundle{}的情形，也不同于稍后 $\mathrm{Ext}$ 的情形中那些相应的命题）。

\emph{证明。}\quad 我们有函子性的\trIsomorphism
\[
\begin{aligned}
\Hom_X(\mathscr{F},\pi^{!?}\omega_Y)\cong{}& \Hom_Y(\pi_*\mathscr{F},\omega_Y) && \text{（$(\pi_*,\pi^{!?})$ 的\trAdjointness）}\\
\cong{}& \mathrm{H}^n(Y,\pi_*\mathscr{F})^{\vee} && \text{（$Y$ 的 $\mathrm{Hom}$ \trSerreDuality）}\\
\cong{}& \mathrm{H}^n(X,\mathscr{F})^{\vee}. && \text{（$\pi$ 的\trAffineness）}
\end{aligned}
\]
\hfill$\square$

你可以把上述论证展开，用 $Y$ 的\trTraceMap{} $t_Y$ 表示出 $X$ 的\trTraceMap{} $t_X$：它就是复合
\[
\mathrm{H}^n(X,\pi^{!?}\omega_Y)\longrightarrow\mathrm{H}^n(Y,\pi_*\pi^{!?}\omega_Y)\longrightarrow\mathrm{H}^n(Y,\omega_Y)\overset{t_Y}{\longrightarrow} k.
\]

\textbf{29.3.12. \trhCorollary。}\ ——\ \emph{$\mathrm{Hom}$ 的\trSerreDuality{}的函子版本与迹版本对一切纯维\trProjective{} $k$-\trScheme{}成立（不带任何 \trCM{}假设）。}

\emph{证明。}\quad 把\mbox{命题~29.3.11} 与\mbox{命题~29.3.5(a)} 以及\mbox{命题~29.1.10} 结合起来。\hfill$\square$

\textbf{29.3.13. \trhProposition。}\ ——\ \emph{若 $X$ 是\trPureDimension{}为 $n$ 的 \trCM{} \trProjective{} $k$-\trScheme，而 $\pi\colon X\to\mathbb{P}^n$ 是一个\trFlat{}\trFinite\trMorphism（如\mbox{命题~29.3.5(b)} 中那样），则 $\mathrm{Ext}$ 的函子版本\trSerreDuality{}对 $X$ 成立，带\trCanonical{}层 $\omega_X:=\pi^{!?}\omega_{\mathbb{P}^n}$。}

与\mbox{命题~29.3.9} 和\mbox{命题~29.3.11} 不同的是，我们并没有对一般的 $\pi\colon X\to Y$ 陈述这条命题；我们将用到「目标是 $\mathbb{P}^n$」这一事实。

\emph{证明。}\quad 论证将与 \S29.2.2 平行：我们将给出诸\trDeltaFunctor{}之间的一个\trIsomorphism
\begin{equation*}\tag{29.3.13.1}\label{eq:ext-serre-duality-iso}
\mathrm{Ext}^i_X(\mathscr{F},\pi^{!?}\omega_{\mathbb{P}^n})\overset{\sim}{\longrightarrow}\mathrm{H}^{n-i}(X,\mathscr{F})^{\vee}.
\end{equation*}
$i=0$ 时的\trIsomorphism{}我们已经有了（\mbox{推论~29.3.12}），所以只须证明 \eqref{eq:ext-serre-duality-iso} 的两侧都是泛\trDeltaFunctor（\S23.2.6）。为此，我们来验证\mbox{定理~23.2.9} 的判据。对 $X$ 上任意\trCoherentSheaf{} $\mathscr{F}$，我们可以找到满态射 $\OO(-m)^{\oplus N}\to\mathscr{F}$（其中 $m\gg 0$，而 $N$ 为某个数），于是只须证明：当 $m\gg 0$ 时，
\[
\mathrm{Ext}^i_X(\OO(-m),\pi^{!?}\omega_{\mathbb{P}^n})=0\qquad\text{与}\qquad \mathrm{H}^{n-i}(X,\OO(-m))^{\vee}=0
\]
对 $i>0$ 与 $m\gg 0$ 成立。对于前者，我们有（当 $i>0$ 且 $m\gg 0$ 时）
\[
\begin{aligned}
\mathrm{Ext}^i_X(\OO(-m),\pi^{!?}\omega_{\mathbb{P}^n})&\cong\mathrm{H}^i\bigl(X,(\pi^{!?}\omega_{\mathbb{P}^n})(m)\bigr) && \text{（\mbox{练习~29.2.K}）}\\
&=0 && \text{（$X$ 的\trSerre{}消失，\mbox{定理~18.1.3(ii)}）。}
\end{aligned}
\]
对于后者，我们有
\[
\begin{aligned}
\mathrm{H}^{n-i}(X,\OO_X(-m))&=\mathrm{H}^{n-i}(X,\OO_X\otimes\pi^*\OO(-m))\\
&\cong\mathrm{H}^{n-i}\bigl(\mathbb{P}^n,(\pi_*\OO_X)\otimes\OO(-m)\bigr)
\end{aligned}
\]
（由\trProjection{}公式，\mbox{练习~18.7.E(b)}）。而 $\pi_*\OO_X$ 是\trLocallyFree{}的（\mbox{练习~24.3.E}），于是由 $\mathbb{P}^n$ 上\trVectorBundle{}的函子版本\trSerreDuality，我们有（当 $m\gg 0$ 时）
\[
\begin{aligned}
\mathrm{H}^{n-i}\bigl(\mathbb{P}^n,(\pi_*\OO_X)\otimes\OO(-m)\bigr)\cong{}& \mathrm{H}^{i}\bigl(\mathbb{P}^n,(\pi_*\OO_X)^{\vee}\otimes\OO(m)\otimes\omega_{\mathbb{P}^n}\bigr)^{\vee}\\
\cong{}& \mathrm{H}^{i}\bigl(\mathbb{P}^n,((\pi_*\OO_X)^{\vee}\otimes\omega_{\mathbb{P}^n})\otimes\OO(m)\bigr)^{\vee}\\
={}& 0 \qquad\text{（$\mathbb{P}^n$ 的\trSerre{}消失，\mbox{定理~18.1.3(ii)}）。}
\end{aligned}
\]
\hfill$\square$

\textbf{29.3.14. \trhCorollary。}\ ——\ \emph{$\mathrm{Ext}$ 的函子版本\trSerreDuality{}对一切纯维的 \trCM{} \trProjective{} $k$-\trScheme{}成立。}

\emph{证明。}\quad 把\mbox{命题~29.3.13} 与\mbox{命题~29.3.5(b)} 以及 \S29.2.2 的射影空间情形结合起来。\hfill$\square$

\textbf{29.3.15.}\ \extrahard\ \emph{\trhRemark：\trVectorBundle{}的\trSerreDuality{}的迹版本。}\quad 下面是一个梗概，说明如何把上述证明加以扩展，以对一切纯维的 \trCM{} \trProjective{} $k$-\trScheme{}证明\trVectorBundle{}的\trSerreDuality{}的迹版本。沿用\mbox{命题~29.3.13} 的记号，\trYoneda{}上积（我们尚未定义它）继以\trTraceMap，便给出一个配对
\[
\mathrm{Ext}^i_X(\mathscr{F},\omega_X)\times\mathrm{H}^{n-i}(X,\mathscr{F})\overset{\cup}{\longrightarrow}\mathrm{H}^n(X,\omega_X)\overset{t}{\longrightarrow} k\,,
\]
它随 $\mathscr{F}$ 函子。这诱导出 \eqref{eq:ext-serre-duality-iso} 的一个函子性的映射。可以证明，这个特定的映射是一个\trDeltaFunctor{}态射。于是由上述论证，这个特定的映射是一个\trIsomorphism。

\textbf{29.3.16.}\quad 正如 \S29.1.13 中所描述的，实际情形还要更好：纯维的\trProper{} $k$-\trScheme{}满足最好的那种\trSerreDuality，即 $\mathrm{Ext}$ 的迹版本。[Stacks, tag~0FVU]

\textbf{29.3.17. 结论。}\quad 本节中的这些证明或许出人意料地容易。$\pi^{!?}$ 一经定义，我们就只用了「它是 $\pi_*$ 的右伴随」这唯一一桩事实。我们付出的代价是：从这个证明里，我们对\trCanonical{}层究竟是什么模样几乎一无所知。例如，若 $X$ \trSmooth，并不能明白 $\omega_X$ 是不是一个\trLineBundle！（我们将在下一节中补救这一点。）

你或许也已经感觉到，这些想法如何通向一个更大的故事：为了在更一般的情形下定义并证明\trDuality，我们正试图找出一个函子 $\pi^!$，使之在我们力所能及的最大普遍性上成为 $\pi_*$ 的\trRightAdjoint。这将引导我们在比\trCoherentSheaf{}更一般的框架中工作——\trCoherentSheaf{}的\trDerivedCategory{}（乃至更远）。
